%MECHANICAL PROBLEMS
%ANIMALS.
%FLYING OBJECTS
%GETTING OFF THE ROAD


\langsection[subsection]{Mechanical Problems} {Неисправности}





\begin{langblock}{0.425}
	\begin{eng}
	
	You can find yourself in an emergency the moment something goes wrong with your motorcycle. In dealing with any mechanical problem, take into account the road and traffic conditions you face. Here are some guidelines that can help you handle mechanical problems safely.
	\end{eng}
	\begin{rus}
		Любая неисправность мотоцикла может мгновенно привести к аварийной ситуации. При устранении любой механической неисправности учитывайте окружающие вас условия. Далее приведены рекомендации, как безопасно действовать при технических неисправностях.
\end{rus}
\end{langblock}

\langsection[subsubsection]{Tire Failure}{Прокол колеса}




\begin{langblock}{0.425}
	\begin{eng}
	
	You will seldom hear a tire go flat. If the motorcycle starts handling differently, it may be a tire failure.  This can be dangerous. You must be able to tell from the way the motorcycle reacts. If one of your tires suddenly loses air, react quickly to keep your balance. Pull off and check the tires.
	
	If the front tire goes flat, the steering will feel heavy. A front-wheel flat is particularly hazardous because it affects your steering. You have to steer well to keep your balance.
	
	If the rear tire goes flat, the back of the motorcycle may jerk or sway from side to side.
	
	If either tire goes flat while riding:
	\begin{itemize}[noitemsep]
		\item Hold handgrips firmly, ease off the throttle, and keep a straight course.
		\item If braking is required, however, gradually apply the brake of the tire that isn't flat, if you are sure which one it is.
		\item When the motorcycle slows, edge to the side of the road, squeeze clutch and stop.
	\end{itemize}
	\end{eng}
	\begin{rus}
		
		Прокол колеса на мотоцикле обычно бесшумный; Вы узнаете о нем по изменению руления, а не по звуку. Если мотоцикл начал вести себя иначе, это может быть из-за прокола.  Это крайне опасно. Вам нужно уметь определять это по реакции мотоцикла. Если одна из шин внезапно спустила, реагируйте быстро, чтобы не потерять контроль. Съедьте на край дороги и проверьте колеса.
		
		Если спустило переднее колесо, руление станет «тяжелым». Прокол переднего колеса особенно опасен, поскольку влияет на управляемость. Чтобы удержать равновесие, нужно четко контролировать руль.
		
		Если спустилась задняя шина, задняя часть мотоцикла может дергаться или раскачиваться из стороны в сторону.
		
		Если во время езды спустилась любая из шин:
		
		\begin{itemize}[noitemsep]
			\item Крепко держитесь за ручки, плавно сбавьте газ и держите прямой курс.
			\item Если требуется торможение, постепенно задействуйте тормоз исправного колеса, конечно, если вы уверены, какого именно.
			\item Когда мотоцикл достаточно замедлится, съезжайте на край дороги, выжмите сцепление и остановитесь.\footnote{Тут речь именно про край дорожного полотна, а не обочину.}
		\end{itemize}
\end{rus}
\end{langblock}


\langsection[subsubsection]{Stuck Throttle}{Заклинило Газ}



\begin{langblock}{0.425}
	\begin{eng}
	
	Twist the throttle back and forth several times. If the throttle cable is stuck, this may free it. If the throttle stays stuck immediately operate the engine cut-off switch and pull in the clutch at the same time. This will remove power from the rear wheel, though engine noise may not immediately decline. Once the motorcycle is under control, pull off and stop.
	
	After you have stopped, check the throttle cable carefully to find the source of the trouble. Make certain the throttle works freely before you start to ride again.
	\end{eng}
	\begin{rus}
		
	Повращайте ручку газа вперед-назад несколько раз. Если трос газа застрял, это может его высвободить. Если ручка газа по-прежнему не слушается,  одновременно выжмите сцепление и кнопку kill switch. Раскрученный мотор затихнет не сразу. Как только мотоцикл успокоится и стабилизируется, съезжайте на обочину и останавливайтесь.
	
	После остановки тщательно осмотрите трос дроссельной заслонки, чтобы найти причину неисправности. Перед тем как продолжить поездку, не спешите, убедитесь, что ручка газа работает как надо.
\end{rus}
\end{langblock}

\langsection[subsubsection]{Wobble}{Шимми, Воблинг}




\begin{langblock}{0.425}
	\begin{eng}
	
	A wobble occurs when the front wheel and handlebars suddenly start to shake from side to side at any speed. Most wobbles can be traced to improper loading, unsuitable accessories, or incorrect tire pressure.  If you are carrying a heavy load, lighten it. If you can't, shift it. Center the weight lower and farther forward on the motorcycle. Make sure tire pressure, spring pre-load, air shocks, and dampers are at the settings recommended for that much weight.  Make sure windshields and fairings are mounted properly.
	
	Check for poorly adjusted steering; worn steering parts; a front wheel that is bent, misaligned, or out of balance; loose wheel bearings or spokes; and swingarm bearings.  If none of these are determined to be the cause, have the motorcycle checked out thoroughly by a qualified professional.
	
	Trying to accelerate out of a wobble will only make the motorcycle more unstable. Instead:
	
	\begin{itemize}[noitemsep]
		\item Grip the handlebars firmly, but don't fight the wobble.
		\item Close the throttle gradually to slow down. Do not apply the brakes; braking could make the wobble worse.
		\item Move your weight as far forward and down as possible.
		\item Pull off the road as soon as you can to fix the problem.
	\end{itemize}
	\end{eng}
	\begin{rus}
		
		«Шимми» (wobble, воблинг) возникает, когда переднее колесо и руль внезапно начинают раскачиваться из стороны в сторону на любой скорости. В большинстве случаев причиной колебаний являются неправильная расстановка груза, неподходящие аксессуары или давление в шинах.  Если вы везете тяжелый груз, уменьшите его вес. Если это невозможно, переместите его. Расположите груз по центру, ниже и ближе к передней части мотоцикла. Убедитесь, что давление в шинах, преднатяг пружин, пневматические амортизаторы и демпферы настроены в соответствии с рекомендациями для данной массы.  Убедитесь, что ветровые стекла и обтекатели установлены правильно.
		
		Проверьте, рулевую колонку; перетянуты или изношены детали рулевого механизма; не погнуто ли, смещено или разбалансировано переднее колесо; не ослаблены ли подшипники колеса или спицы; а также подшипники маятника.  Если вы не нашли причину, то поручите квалифицированному специалисту провести тщательную проверку.
		
		Попытка «выйти из колебаний через разгон» только сделает хуже. Вместо этого:\footnote{Степень расколбаса бывает разной: от удивления до паники. Чаще случается на спорт-байках. Если вы справились -- считайте вам крупно повезло.}
		
		\begin{itemize} [noitemsep]
			\item Крепко держите руль, но не сопротивляйтесь колебаниям.
			\item Плавно сбросьте газ, чтобы снизить скорость. Не нажимайте на тормоза; торможение может усугубить колебания.
			\item Перенесите вес тела вперед и вниз на сколько это возможно.
			\item Как только станет возможным, съезжайте с дороги и ищите проблему.
		\end{itemize}
\end{rus}
\end{langblock}



\langsection[subsubsection]{Drivetrain Problems}{Проблемы с трансмиссией}


\begin{langblock}{0.425}
	\begin{eng}
	
	A chain that slips or breaks while you're riding could lock the rear wheel and cause your cycle to skid.  Chain slippage or breakage can be avoided by proper maintenance.
	
	\begin{description}
		\item[Slippage] If the chain slips when you try to speed up quickly or ride uphill, pull off the road. Check the chain and sprockets. Tightening the chain may help. If the problem is a worn or stretched chain or worn or bent sprockets, replace the chain, the sprockets, or both before riding again.
		\item[Breakage] You'll notice an instant loss of power to the rear wheel.  Close the throttle and brake to a stop.
	\end{description}
	\end{eng}
	\begin{rus}
		
		Цепь, перескакивающая или рвущаяся в движении, может заблокировать заднее колесо и вызвать занос мотоцикла. Проскальзывания или обрыва цепи можно избежать, если делать тех. обслуживание вовремя.
		
		\begin{description}[noitemsep]
			\item[Проскальзывание] Если цепь проскальзывает при попытке резкого ускорения или при подъеме в гору, съезжайте с дороги. Проверьте цепь и звездочки. Помочь может натяжение цепи. Если проблема заключается в изношенной или растянутой цепи либо в изношенных или погнутых звездочках, замените цепь, звездочки или и то, и другое, прежде чем снова садиться на мотоцикл.
			\item[Разрыв] Вы заметите мгновенную потерю тяги на заднем колесе.  Сбросьте газ и затормозите до полной остановки.
		\end{description}
\end{rus}
\end{langblock}

\langsection[subsubsection]{Engine Seizure}{Клин двигателя}


\begin{langblock}{0.425}
	\begin{eng}
	
	When the engine locks or freezes it is usually low on oil. The engine's moving parts can't move smoothly against each other, and the engine overheats. The first sign may be a loss of engine power or a change in the engine's sound. Squeeze the clutch lever to disengage the engine from the rear wheel. Pull off the road and stop. Check the oil. If needed, oil should be added as soon as possible or the engine will seize. When this happens, the effect is the same as a locked rear wheel. Let the engine cool before restarting.
	\end{eng}
	\begin{rus}
		
		Когда двигатель «заклинивает» или он «залипает», это, как правило, связано с низким уровнем масла. Движущиеся детали не могут плавно скользить друг по другу, в результате чего двигатель перегревается. Первым признаком может стать потеря мощности или изменение звука. Выжмите сцепление, чтобы отсоединить двигатель от заднего колеса. Съезжайте с дороги и остановитесь. Проверьте уровень масла. При необходимости масло следует долить как можно скорее, иначе двигатель заклинит. Если это произойдет, последствия будут такими же, как при заблокированном заднем колесе. Перед повторным запуском дайте двигателю остыть.
\end{rus}
\end{langblock}

\langsection[subsection]{Animals}{Животные}



\begin{langblock}{0.425}
	\begin{eng}
	
	Naturally, you should do everything you safely can to avoid hitting an animal. If you are in traffic, however, remain in your lane. Hitting something small is less dangerous to you than hitting something big -- like a car.
	
	Motorcycles seem to attract dogs. If you are chased, downshift and approach the animal slowly. As you approach it, accelerate away and leave the animal behind. Don't kick at an animal. Keep control of your motorcycle, and look to where you want to go.
	
	For larger animals (deer, elk, cattle) brake and prepare to stop, they are unpredictable.
	\end{eng}
	\begin{rus}
		
		Разумеется, вы должны сделать все, чтобы по возможности без риска для себя избежать наезда на животное. Однако, если вы находитесь в потоке транспорта, оставайтесь на своей полосе движения. Наезд на небольшое животное менее опасен для вас, чем столкновение с чтем-то крупным "--- например, автомобилем.
		
		Мотоциклы часто притягивают собак. Если вас преследуют, переключитесь на пониженную и медленно приближайтесь к животному. Приблизившись, ускорьтесь и уезжайте прочь. Не пытайтесь бить ногой по животному. Сохраняйте контроль над мотоциклом и смотрите в том направлении, куда хотите ехать.
		
		В случае с крупными животными (олени, лоси, крупный рогатый скот) притормозите и будьте готовы остановиться "--- их поведение непредсказуемо.
\end{rus}
\end{langblock}

\langsection[subsection]{Flying Objects}{Все что летает}


\begin{langblock}{0.425}
	\begin{eng}
	
	From time to time riders are struck by insects, cigarettes thrown from cars, or pebbles kicked up by the tires of the vehicle ahead. If you are wearing face protection, it might get smeared or cracked, making it difficult to see. Without face protection, an object could hit you in the eye, face, or mouth. Whatever happens, keep your eyes on the road and your hands on the handlebars. When safe, pull off the road and repair the damage.
	\end{eng}
	\begin{rus}
		
		Время от времени вы столкнетесь с насекомыми, сигаретами, выброшенными из окна, или гравием из-под колес впереди едущих авто. Без защитной маски они могли попасть вам в глаз, лицо или рот. Если вы носите защитную маску, на ней могут появиться разводы или трещины, затрудняющие обзор. Что бы ни случилось, не отрывайте взгляд от дороги и держите руки на руле. Когда это будет безопасно, съедьте с дороги и устраните последствия.
\end{rus}
\end{langblock}


\langsection[subsection]{Getting off the Road}{Съезд на обочину}


\begin{langblock}{0.425}
	\begin{eng}
	
		If you need to leave the road to check the motorcycle (or just to rest for a while), be sure you:
		
		\begin{itemize}[noitemsep]
			\item Check the roadside -- Make sure the surface of the roadside is firm enough to ride on. If it is soft grass, loose sand, or if you're just not sure about it, slow way down before you turn onto it.
			\item Signal -- Drivers behind might not expect you to slow down. Give a clear signal that you will be slowing down and changing direction. Check your mirror and make a head check before you take any action.
			\item Pull off the road -- Get as far off the road as you can. It can be very hard to spot a motorcycle by the side of the road. You don't want someone else pulling off at the same place you are.
			\item Park carefully -- Loose and sloped shoulders can make setting the side or center stand difficult.
		\end{itemize}
	\end{eng}
	\begin{rus}
		
		Если вам нужно съехать с дороги, чтобы проверить мотоцикл (или просто немного отдохнуть), обязательно:
		
		\begin{itemize}[noitemsep]
			\item Оцените состояние обочины "--- убедитесь, что поверхность обочины достаточно твердая. Если это трава, песок или вы просто не уверены, значительно снизьте скорость, прежде чем съезжать на обочину.
			\item Подайте сигнал поворота "--- Водители, едущие сзади, могут не ожидать, что вы замедлитесь. Четко дайте понять, что собираетесь снизить скорость и изменить направление движения. Перед тем как предпринять какие-либо действия, посмотрите в зеркало и оглянитесь назад.
			\item Съезжайте с дороги "--- Отъедьте от дороги как можно дальше. Мотоцикл, стоящий на обочине очень трудно заметить. Вы же не хотите, чтобы кто-то еще съехал на то месте, где вы стоите.
			\item Осторожно паркуйтесь "--- мягкая или с уклоном обочина может затруднить постановку на боковую или центральную подножку.
		\end{itemize}
\end{rus}
\end{langblock}









